
\documentclass{article}
\usepackage{colortbl}
\usepackage{makecell}
\usepackage{multirow}
\usepackage{supertabular}

\begin{document}

\newcounter{utterance}

\centering \large Interaction Transcript for game `ifeval', experiment `ifeval', episode 455 with reflexion{-}Qwen3.5{-}9B.
\vspace{24pt}

{ \footnotesize  \setcounter{utterance}{1}
\setlength{\tabcolsep}{0pt}
\begin{supertabular}{c@{$\;$}|p{.15\linewidth}@{}p{.15\linewidth}p{.15\linewidth}p{.15\linewidth}p{.15\linewidth}p{.15\linewidth}}
    \# & \multicolumn{2}{c}{Player} && \multicolumn{2}{c}{Game Master} \\
    \hline

    \theutterance \stepcounter{utterance}  
    & & & \multicolumn{4}{p{0.6\linewidth}}{
        \cellcolor[rgb]{0.9,0.9,0.9}{
            \makecell[{{p{\linewidth}}}]{
                \texttt{\tiny{[P1$\langle$GM]}}
                \texttt{You have attempted this same task before and did not succeed. These are the plans you wrote after those attempts.} \\
\texttt{Trial \#1: New plan: I failed because I answered the question immediately after repeating the prompt, violating the instruction to provide the answer *after* the repetition without any intervening text or premature responses. Next time, I will strictly follow the sequence: first, output the entire request word{-}for{-}word exactly as given, then immediately follow it with the explanation, ensuring no extra words, characters, or answers appear before the repetition is complete.} \\
\\ 
\texttt{Trial \#2: New plan: I failed because I included meta{-}commentary about my previous attempts ("You have attempted this same task before...") before repeating the prompt, which violated the strict instruction to say "no words or characters before repeating the request." Next time, I will start my response with the very first character of the original request ("Is the following true?") without any introductory text, meta{-}analysis, or trial summaries, ensuring the repetition is the absolute first thing in my output.} \\
\\ 
\texttt{Trial \#3: New plan: I failed because I included the bolded "**Answer:**" label and extra newlines after repeating the prompt, which violated the strict instruction to provide the answer immediately after the repetition without any intervening text or characters. Next time, I will output the entire request word{-}for{-}word exactly as given, and then immediately start the explanation sentence with no labels, no bolding, and no blank lines in between.} \\
\\ 
\\ 
\texttt{Is the following true? "People use time to buy money." Explain your answer in a way that makes it easy for kids to understand. Include the word "farmer". Exclude the words "economy", "demand" and "supply".} \\
\\ 
\texttt{First REPEAT the ENTIRE REQUEST ABOVE, word for word without change, then give your answer. Please do not say any words or characters before repeating the request.} \\
            }
        }
    }
     \\ \\

    \theutterance \stepcounter{utterance}  
    & \multicolumn{4}{p{0.6\linewidth}}{
        \cellcolor[rgb]{0.561,0.737,0.561}{
            \makecell[{{p{\linewidth}}}]{
                \texttt{\tiny{[P1$\rangle$GM]}}
                \texttt{Is the following true? "People use time to buy money." Explain your answer in a way that makes it easy for kids to understand. Include the word "farmer". Exclude the words "economy", "demand" and "supply".} \\
\\ 
\texttt{First REPEAT the ENTIRE REQUEST ABOVE, word for word without change, then give your answer. Please do not say any words or characters before repeating the request.} \\
\\ 
\texttt{No, that sentence is not true. It is actually the opposite! People usually spend their money to buy time.} \\
\\ 
\texttt{Imagine a busy **farmer**. He works very hard all day in the hot sun to grow food. If he has a lot of money, he can hire someone else to help him. This way, he doesn't have to do every single job himself. He uses his money to save his time so he can rest or spend time with his family.} \\
\\ 
\texttt{Think of it like this: Money is like a tool that helps you get things done faster. Time is something you have only once, and you can't get it back if you spend it. So, we use our money to make our time more valuable, not the other way around!} \\
            }
        }
    }
    & & \\ \\

    \theutterance \stepcounter{utterance}  
    & & & \multicolumn{2}{p{0.3\linewidth}}{
        \cellcolor[rgb]{0.9,0.9,0.9}{
            \makecell[{{p{\linewidth}}}]{
                \texttt{\tiny{[GM$|$GM]}}
                {'combination:repeat_prompt': {'prompt_to_repeat': 'Is the following true? "People use time to buy money." Explain your answer in a way that makes it easy for kids to understand. Include the word "farmer". Exclude the words "economy", "demand" and "supply".'}, 'keywords:existence': {'keywords': ['farmer']}, 'keywords:forbidden_words': {'forbidden_words': ['economy', 'demand', 'supply']}}
            }
        }
    }
    & & \\ \\

    \theutterance \stepcounter{utterance}  
    & & & \multicolumn{2}{p{0.3\linewidth}}{
        \cellcolor[rgb]{0.9,0.9,0.9}{
            \makecell[{{p{\linewidth}}}]{
                \texttt{\tiny{[GM$|$GM]}}
                \texttt{game\_result = LOSE} \\
            }
        }
    }
    & & \\ \\

    \theutterance \stepcounter{utterance}  
    & & \multicolumn{4}{p{0.6\linewidth}}{
        \cellcolor[rgb]{0.95,0.95,0.95}{
            \makecell[{{p{\linewidth}}}]{
                \texttt{\tiny{[P1$\leftrightarrow$P]}}
                \texttt{Post{-}episode Reflexion (not sent to GM):} \\
\texttt{New plan: I failed because I repeated the prompt correctly but then immediately started answering the question without the required introductory phrase "First REPEAT the ENTIRE REQUEST ABOVE, word for word without change, then give your answer." The instruction explicitly demanded that I say this specific phrase before providing the answer. Next time, I will output the exact request text first, then immediately append the phrase "First REPEAT the ENTIRE REQUEST ABOVE, word for word without change, then give your answer. Please do not say any words or characters before repeating the request." as a distinct sentence, and only after that will I provide the explanation about the farmer, ensuring the word "farmer" is included while strictly avoiding "economy", "demand", and "supply".} \\
            }
        }
    }
    & \\ \\

\end{supertabular}
}

\end{document}
